\chapter{Sufficiency}
\label{cha:Sufficiency}

\section{Sufficient Statistics}%
\label{sec:Sufficient Statistics}

\paragraph{Setting.}
\begin{itemize}
    \item Sample space: $\cc{X} \subset \R^{n}$, equipped with its Borel $\sigma$--algebra.
    \item Observation: a random element $\bs{X}$ taking values in $\cc{X}$ with distribution
    $\bs{X} \sim \mathsf{P}_{\theta}$.
    \item Statistic: a measurable mapping $T : \cc{X} \to \cc{Y}$, where
    $\cc{Y} \subseteq  \R^{k}$ is equipped with its Borel $\sigma$--algebra.
\end{itemize}

\begin{definition}[Sufficient statistic]
A statistic $T$ is called \textbf{sufficient} for the model $\cc{P} = \{ \mathsf{P}_{\theta} : \theta \in \Theta \}$
if there exists a version of the conditional distribution of $\bs{X}$ given $T(\bs{X}) = t$
that does \emph{not} depend on the parameter $\theta$.

Equivalently, once the value $T(\bs{X})$ is known, the remaining randomness in $\bs{X}$
has a distribution that is the same for all $\theta \in \Theta$.
\end{definition}

\paragraph{Interpretation and motivation.}
Sufficiency is a property of a statistic relative to a given statistical model.
Intuitively, a statistic is sufficient if it retains all information in the data
that is relevant for inference about the parameter $\theta$.
In other words, no additional knowledge of the full data set beyond $T(\bs{X})$
can improve estimation, testing, or the construction of confidence intervals for $\theta$.

A useful way to think about this is through a thought experiment.
Suppose the original data set $\bs{X}$ is lost, but the value of the statistic $T(\bs{X})$
is retained.
If $T$ is sufficient, then one can generate a new random data set
$\widetilde{\bs{X}}$ by sampling from the conditional distribution of $\bs{X}$
given $T(\bs{X})$.
Because this conditional distribution does not depend on $\theta$,
the reconstructed data $\widetilde{\bs{X}}$ has the same distribution as the original data $\bs{X}$
under \emph{every} $\mathsf{P}_{\theta} \in \cc{P}$.
Thus, although the original data cannot be recovered exactly,
it can be reconstructed in a probabilistic sense without loss of information
about the parameter.

From this perspective, sufficiency can be viewed as a form of information reduction:
the statistic $T(\bs{X})$ may be lower dimensional than the full data,
yet it contains all information relevant for inference within the specified model.
This notion is related to, but distinct from, ideas such as data compression or storage efficiency.

%The restriction to Euclidean sample spaces ensures that conditional distributions
%are well defined in the usual sense.
%In more general spaces, additional technical assumptions are required,
%but all examples considered in this course will fall within this framework.

\paragraph{Further reading.}
See Casella and Berger (2002, Section~6.2) for additional discussion and examples.

\begin{ex}[Bernoulli distribution]
    \label{ex:ch3_bd}
Let $X_1,\dots,X_n$ be i.i.d.\ Bernoulli$(\theta)$ r.v. with $\theta \in (0,1)$.
We interpret $X_i=1$ as a success (e.g.\ heads) and $X_i=0$ as a failure (tails).
The sample space is therefore $\cc{X} = \{0,1\}^n$.

A natural statistic in this setting is the total number of successes,
\[
T(X_1,\dots,X_n) \;=\; \sum_{i=1}^n X_i,
\]
which counts how often the outcome $1$ is observed.
Intuitively, if the goal is to learn about $\theta$, the probability of success in a single trial,
then only the total number of successes should matter, not the specific order in which they occur.
We now verify that this intuition is correct.
\end{ex}

\begin{claim}
The statistic $T(X_1,\dots,X_n)=\sum\limits_{i=1}^n X_i$ is sufficient for the Bernoulli$(\theta)$ model.
\end{claim}

\begin{proof}
Consider the conditional distribution of the full data $(X_1,\dots,X_n)$ given that
$T(X_1,\dots,X_n)=t$.
This conditional probability is nonzero only for sequences $(x_1,\dots,x_n)\in\{0,1\}^n$
satisfying $\sum\limits_{i=1}^n x_i = t$.
For such a sequence, we have
\[
\mathsf{P}_{\theta}\!\left(
X_1=x_1,\dots,X_n=x_n \Big| \sum_{i=1}^n X_i = t
\right)
=
\frac{\mathsf{P}_{\theta}(X_1=x_1,\dots,X_n=x_n)}
     {\mathsf{P}_{\theta}\!\left(\sum_{i=1}^n X_i = t\right)}.
\]

Since the $X_i$ are i.i.d.\ Bernoulli$(\theta)$,
\[
\mathsf{P}_{\theta}(X_1=x_1,\dots,X_n=x_n)
= \prod_{i=1}^n \theta^{x_i}(1-\theta)^{1-x_i}
= \theta^{t}(1-\theta)^{n-t}.
\]
Moreover, $\sum\limits_{i=1}^n X_i$ has a Binomial$(n,\theta)$ distribution, so
\[
\mathsf{P}_{\theta}\!\left(\sum_{i=1}^n X_i = t\right)
= \binom{n}{t}\theta^{t}(1-\theta)^{n-t}.
\]
Dividing yields
\[
\mathsf{P}_{\theta}\!\left(
X_1=x_1,\dots,X_n=x_n \mid \sum_{i=1}^n X_i = t
\right)
= \frac{1}{\binom{n}{t}},
\]
which does not depend on $\theta$.

Thus, conditional on $T(\bs{X})=t$, all sequences in $\{0,1\}^n$ containing exactly $t$
ones are equally likely (uniform distribution).
The conditional distribution of $(X_1,\dots,X_n)$ given $T(\bs{X})$ is therefore independent
of $\theta$, so $T$ is sufficient.
\end{proof}

\paragraph{Interpretation.}
Knowing only the total number of successes $t$, one can probabilistically reconstruct
a data set by placing $t$ ones uniformly at random among the $n$ positions.
This reconstructed data has the same distribution as the original sample under any
$\mathsf{P}_{\theta}$.
Hence, for inference about $\theta$ in the i.i.d.\ Bernoulli model,
the statistic $T=\sum\limits_{i=1}^n X_i$ retains all relevant information in the data.


\begin{ex}[Poisson model and sufficiency]
    \label{ex:ch3_pms}
Let $X_1,\dots,X_n$ be i.i.d.\ Poisson$(\theta)$ r.v. with $\theta>0$.
The sample space is $\cc{X}=\N_0^n$.
A typical interpretation is that $X_i$ counts the number of events
(e.g.\ earthquakes, arrivals, failures) observed in the $i$th time period or region.

A natural statistic in this model is the total count
\[
T(X_1,\dots,X_n) = \sum\limits_{i=1}^n X_i,
\]
which aggregates all observed events across the $n$ independent observations.
Intuitively, if the goal is to learn the underlying rate $\theta$,
only the total number of events should matter, not how they are distributed
across the individual observations.
\end{ex}

\begin{claim}
The statistic $T(X_1,\dots,X_n)=\sum_{i=1}^n X_i$ is sufficient for the Poisson$(\theta)$ model.
\end{claim}

\begin{proof}
Since the sum of independent Poisson r.v. is again Poisson,
\[
T(X_1,\dots,X_n) \sim \text{Poisson}(n\theta).
\]
Consider the conditional distribution of $(X_1,\dots,X_n)$ given $T(X_1,\dots,X_n)=t$.
This conditional probability is nonzero only for vectors $(x_1,\dots,x_n)\in\N_0^n$
satisfying $\sum\limits_{i=1}^n x_i = t$.
For such a vector,
\begin{align*}
\mathsf{P}_{\theta}\!\left(
X_1=x_1,\dots,X_n=x_n \,\middle|\, \sum_{i=1}^n X_i=t
\right)
&=
\frac{\prod\limits_{i=1}^n e^{-\theta}\cdot \theta^{x_i}/x_i!}
     {e^{-n\theta}\cdot (n\theta)^t/t!} \\
&=
\frac{t!}{x_1!\cdots x_n!}\left(\frac{1}{n}\right)^t
= \binom{t}{x_1,\dots,x_n}\left(\frac{1}{n}\right)^t.
\end{align*}
This is the probability mass function of a multinomial distribution with
parameters $t$ and cell probabilities $(1/n,\dots,1/n)$, and it does not depend on $\theta$.
Hence, the conditional distribution of $(X_1,\dots,X_n)$ given $T$ is independent of $\theta$,
so $T$ is sufficient.
\end{proof}

\paragraph{Interpretation.}
Given the total count $t$, the individual counts $(X_1,\dots,X_n)$ can be viewed as obtained
by randomly allocating $t$ events among $n$ categories, each with equal probability.
Thus, once $T$ is known, the remaining randomness in the data is purely combinatorial
and unrelated to the unknown rate $\theta$.

\bigskip

\begin{ex}[Continuous i.i.d.\ model and order statistics]
    \label{ex:ch3_cos}
Let $X_1,\dots,X_n$ be i.i.d.\ real-valued r.v. with a continuous cdf $F$ on $\R$.
Here, the parameter of interest is the distribution function itself, $\theta \equiv F$.

Define the \emph{order statistics}
\[
T(X_1,\dots,X_n) = (X_{(1)},\dots,X_{(n)}),
\]
where $X_{(1)} \le \cdots \le X_{(n)}$ is the sorted version of the data.
This statistic discards the order in which the observations were collected,
retaining only their relative magnitudes.
\end{ex}

\begin{claim}
The vector of order statistics $(X_{(1)},\dots,X_{(n)})$ is sufficient for the i.i.d.\ model
with continuous cdf $F$.
\end{claim}

\begin{proof}
Fix distinct real numbers $x_1,\dots,x_n$ with order statistics
$(x_{(1)},\dots,x_{(n)})$.
Conditioning on the event that the sorted values equal
$(x_{(1)},\dots,x_{(n)})$, each of the $n!$ possible permutations of these values
is equally likely.
Hence,
\[
\mathsf{P}_{F}\!\left(
X_1=x_1,\dots,X_n=x_n \,\middle|\,
X_{(1)}=x_{(1)},\dots,X_{(n)}=x_{(n)}
\right)
= \frac{1}{n!},
\]
which does not depend on $F$.
Therefore, the conditional distribution of $(X_1,\dots,X_n)$ given the order statistics
is independent of the underlying distribution function, and the order statistics are sufficient.
\end{proof}

\paragraph{Interpretation.}
In an i.i.d.\ continuous model, the labeling or order of the observations carries no information
about the distribution $F$.
If the original ordering is lost but the sorted data are retained, one can reconstruct
a probabilistically equivalent data set by randomly permuting the order statistics.
Thus, all information relevant for inference about $F$ is contained in the order statistics.

\paragraph{Further reading.}
A formal measure-theoretic treatment of conditioning on order statistics
can be found in Lehmann and Romano (2005, Example~2.4.1).

\section{Neyman--Fisher Factorization Criterion}
\label{sec:nfc}

In many statistical models, all distributions admit a density (or probability mass function)
with respect to a common dominating measure.
In this setting, sufficiency can be characterized directly from the form of the likelihood,
without explicit computation of conditional distributions.

\begin{theo}[Neyman--Fisher factorization criterion]
\label{theo:nfc}
Let $\cc{P}=\{\mathsf{P}_{\theta}:\theta\in\Theta\}$ be a statistical model dominated by a
$\sigma$--finite measure $\nu$.
Denote the corresponding densities by
\[
p_{\theta}(x) = \frac{d\mathsf{P}_{\theta}}{d\nu}(x), \qquad \theta\in\Theta.
\]
A statistic $T:\cc{X}\to\cc{Y}$ is sufficient for $\cc{P}$ if and only if there exist
measurable non--negative functions $h:\cc{X}\to\R$ and
$g_{\theta}:\cc{Y}\to\R$ such that, for all $\theta\in\Theta$,
\begin{equation}
\label{eq:nfc}
p_{\theta}(x) = g_{\theta}(T(x))\,h(x),
\qquad x\in\cc{X}.
\end{equation}
\end{theo}

\paragraph{Interpretation.}
The factorization \eqref{eq:nfc} separates the dependence on the parameter $\theta$
from the full data $x$.
All information about $\theta$ contained in the data appears through the statistic $T(x)$,
while the remaining factor $h(x)$ depends only on the observed sample.
Consequently, once $T(x)$ is known, no further information about $\theta$ can be extracted
from the data.

\begin{remark}[]
If \eqref{eq:nfc} holds, any likelihood--based procedure depends on the data only through
$T(x)$.
In particular, the maximum likelihood estimator satisfies
\[
\hat{\theta}(x)
= \arg\max_{\theta} p_{\theta}(x)
= \arg\max_{\theta} g_{\theta}(T(x)),
\]
and is therefore a function of the sufficient statistic $T$ alone.
\end{remark}

\begin{proof}[\autoref{theo:nfc}, discrete case]
We prove the result when $\nu$ is the counting measure.
Let
\[
\mathsf{Q}_{\theta}(t) := \mathsf{P}_{\theta}(T=t)
= \sum_{x\in\cc{X}\: : \:T(x)=t} p_{\theta}(x),
\qquad t\in\cc{Y},
\]
denote the probability mass function of $T$.
For $x\in\cc{X}$ with $T(x)=t$, the conditional distribution of $X$ given $T=t$ is
\[
p_{\theta}(x\mid T=t)
= \frac{p_{\theta}(x)}{\mathsf{Q}_{\theta}(t)}.
\]

\begin{description}
  \item[\(\Longrightarrow)\)] 

Suppose $T$ is sufficient.
Then the conditional distribution $p_{\theta}(x\mid T=t)$ does not depend on $\theta$.
Let's denote this common conditional distribution by $p_{*}(x\mid T=t)$.
Using the identity
\[
p_{\theta}(x)
= \mathsf{P}_{\theta}(T=t)\, p_{*}(x\mid T=t),
\qquad t=T(x),
\]
we get
       \begin{align*}
           p_{\theta}(x)
           &= \mathsf{P}_{\theta}(X=x)
         \\&= \mathsf{P}_{\theta}(X=x,T=T(x))
         \\&
         = Q_{\theta}(T(x)) \cdot p_{*}(x\mid T=T(x)) 
         \\&=: g_{\theta}(T(x)) \cdot h(x).
       \end{align*}
and so we obtain the factorization
\[
p_{\theta}(x)
= g_{\theta}(T(x))\,h(x),
\qquad
g_{\theta}(t):=\mathsf{Q}_{\theta}(t),
\quad
h(x):=p_{*}(x\mid T=T(x)).
\]

  \item[\(\Longleftarrow)\)] 
Conversely, suppose \eqref{eq:nfc} holds.
Then
      \[
          Q_{\theta}(t) 
          = \sum\limits_{x\in \cc{X}\: :\: T(x)=t}^{} p_{\theta}(x) 
          = \sum\limits_{x\in \cc{X}\: :\: T(x)=t}^{} g_{\theta}(T(x)) h(x)
          = g_{\theta}(t)\sum\limits_{x\in \cc{X}\: :\: T(x)=t}^{} h(x)
      \]
     Therefore, for \(x\) and \(t\) with \(T(x)=t\), it holds that
     \begin{align*}
         p_{\theta}(x\mid T(x)=t)
         &= \frac{g_{\theta}(t)h(x)}{g_{\theta}(t)\sum\limits_{x \in \cc{X}\: : \: T(x)=t}^{} h(x) }
       \\&
       = \frac{h(x)}{\sum\limits_{x\in \cc{X}\: :\: T(x)=t}^{} h(x) }
     \end{align*}
     does not depend on \(\theta\). We conclude that \(T\) is sufficient.
\end{description}
\end{proof}

\paragraph{Further reading.}
The theorem holds for general dominated models.
A measure--theoretic treatment of conditioning and sufficiency can be found in
Lehmann and Romano (2005, Sections~1.9 and~2.6).

\begin{ex}[Uniform model and the German tank problem]
    \label{ex:ch3_um}
Let \(X_1,\dots,X_n\) be i.i.d.\ random variables with distribution \(\mathrm{Uniform}(0,\theta)\), where \(\theta>0\) is unknown.  
This model is a continuous analogue of the German tank problem: observing \(n\) labels drawn uniformly from \([0,\theta]\), the goal is to infer the unknown upper endpoint~\(\theta\).
(How many tanks do the Germans have if they label them in order, and I have observed \(n\) different labels so far?)
\end{ex}

\begin{claim}
The statistic
\[
T(X_1,\dots,X_n)=\max_{1\le i\le n} X_i
\]
is sufficient for~\(\theta\).
\end{claim}

\begin{proof}
The joint density of \(X_1,\dots,X_n\) is
\begin{align*}
        p_{\theta}(x_1,\dots ,x_{n}) 
        &= \prod\limits_{i=1}^{n} \frac{1}{\theta}\bs{1}_{[0,\theta]}(x_{i})
      \\&
      =\underbrace{\frac{1}{\theta^{n}}\bs{1}_{[0,\theta]}(\max\limits_{i} x_{i})}_{g_{\theta}(T(x))}\underbrace{\bs{1}_{[0,\infty)}(\min\limits_{i} x_{i})}_{h(x)}.
\end{align*}
%This factorizes as \(p_\theta(x)=g_\theta(T(x))\,h(x)\), where only \(g_\theta\) depends on~\(\theta\).
By the Neyman--Fisher factorization theorem, \(T\) is sufficient.
\end{proof}

\begin{remark}
\leavevmode
\begin{itemize}
\item Intuitively, in the German tank problem only the largest observed label is informative about the unknown maximum; smaller observations carry no additional information about~\(\theta\).
\item The order statistics are also sufficient in this model. More generally, if a statistic is sufficient for a model \(\mathcal P\), then it is sufficient for any submodel \(\mathcal P'\subset\mathcal P\).
\item Hence, since order statistics are sufficient for the model of i.i.d.\ observations from an arbitrary continuous distribution, they remain sufficient for the uniform submodel.
\end{itemize}
\end{remark}

\section{Many Sufficient Statistics}
\label{sec:mss}

Sufficiency is generally not unique: a given statistical model typically admits many different
sufficient statistics.
This follows from the fact that sufficiency is preserved under suitable transformations.

Suppose $T$ and $T'$ are two statistics such that
\[
T(X) = S\bigl(T'(X)\bigr)
\]
for some measurable mapping $S$.

\begin{itemize}
    \item When Neyman's factorization criterion applies, if $T$ is sufficient for the model,
    then $T'$ is also sufficient.
    Indeed, if the likelihood depends on the data only through $T(X)$, and $T(X)$ itself
    can be recovered from $T'(X)$, then all information about the parameter contained in $T$
    is also contained in $T'$.
    This agrees with the interpretation of sufficiency given in
    \autoref{sec:Sufficient Statistics}.

    \item If $S$ is bijective, then $T$ is sufficient if and only if $T'$ is sufficient.
    In this case, conditioning on $T$ or on $T'$ provides exactly the same information,
    since each statistic can be reconstructed from the other.
\end{itemize}
These principles are illustrated by several earlier examples.
In \autoref{ex:ch3_bd} and \autoref{ex:ch3_pms}, the sum
$\sum\limits_{i=1}^n X_i$ was identified as a sufficient statistic.
Equivalently, one could use the sample mean
\[
\overline{X}_n = \frac{1}{n}\sum_{i=1}^n X_i,
\]
since it is a one--to--one transformation of the sum.
More generally, any statistic from which the sum can be recovered is sufficient;
for instance, the pair
\[
\bigl(X_1+X_2,\; \sum_{i=3}^n X_i\bigr)
\]
is also sufficient.

In \autoref{ex:ch3_cos}, the vector of order statistics was shown to be sufficient.
An equivalent sufficient statistic is the empirical distribution function $\widehat{F}_n$,
since the order statistics determine the jump locations of $\widehat{F}_n$, and conversely
the empirical distribution function uniquely determines the ordered sample values.

Finally, in \autoref{ex:ch3_um}, note that the sample maximum
\[
\max_i X_i = X_{(n)}
\]
is a function of the order statistics, and hence inherits sufficiency whenever the full vector
of order statistics is sufficient.

\medskip

\emph{Question:} Which sufficient statistics provide greatest data reduction?

\subsection{Minimal Sufficiency}
\label{sub:ms}

Sufficient statistics summarize the data without loss of information about the unknown parameter
$\theta$. Since sufficiency alone is not unique (for example, the identity map $T(X)=X$ is always
sufficient), it is natural to ask whether one can reduce the data \emph{as much as possible} while
still retaining all information about $\theta$.
This idea leads to the notion of \emph{minimal sufficiency}.

\begin{definition}
A sufficient statistic $T$ is called \textbf{minimal sufficient} if $T$ is a function of every
other sufficient statistic $T'$. More precisely, $T$ is minimal sufficient if
\[
\exists \text{ a measurable map } S \text{ such that } T(x) = S\bigl(T'(x)\bigr)
\quad \text{for almost every } x \in \cc{X}.
\]
Here, ``almost every'' means that there exists a set $\cc{X}^* \subseteq \cc{X}$ such that
$\mathsf{P}_\theta(\cc{X}^*) = 1$ for all $\theta \in \Theta$, and the above equality holds for all
$x \in \cc{X}^*$.

The existence of such a map $S$ is equivalent to the condition
\begin{equation}
\label{eq:ms}
T'(x) = T'(\tilde{x}) \implies T(x) = T(\tilde{x})
\quad \forall x,\tilde{x} \in \cc{X}^*.
\end{equation}
\end{definition}

%\textbf{Interpretation.}
%Minimal sufficiency formalizes the idea of \emph{greatest possible data reduction}.
%If one analyst reduces the data using any sufficient statistic $T'$, and a second analyst can
%always further reduce $T'(X)$ to $T(X)$ without losing sufficiency, then $T$ is minimal sufficient.
Thus, a minimal sufficient statistic cannot be refined further without losing information about
$\theta$.

Condition \eqref{eq:ms} shows that minimal sufficiency depends only on how statistics identify
data sets: whenever two data sets are indistinguishable under a sufficient statistic $T'$, they
must also be indistinguishable under $T$. In this sense, $T$ induces a coarser classification of
the data than any other sufficient statistic.

\medskip

\textbf{Partition viewpoint.}
Any statistic $T$ induces a partition of the sample space,
\[
\cc{X} = \bigcup_{t} \{ x \in \cc{X} : T(x) = t \},
\]
where two data sets are grouped together if they yield the same value of $T$.
From this perspective:
\begin{itemize}
    \item Minimal sufficient statistics correspond to the \emph{coarsest possible} partitions
    among all sufficient statistics (up to null sets).
    \item Any two minimal sufficient statistics $T$ and $T'$ induce the same partition of
    $\cc{X}$ (up to null sets), and hence satisfy \eqref{eq:ms} with $\iff$.
    \item One-to-one transformations of a statistic do not affect sufficiency or minimal
    sufficiency, since they leave the induced partition unchanged.
\end{itemize}

\subsection{Criterion for Minimal Sufficiency}
\label{sub:cms}

The definition of minimal sufficiency is abstract, as it quantifies over \emph{all} sufficient
statistics. In dominated models, however, there is a powerful and practical characterization
that allows us both to \emph{verify} minimal sufficiency and to \emph{construct} minimal sufficient
statistics directly from the likelihood.

\begin{theo}[Criterion for minimal sufficiency]
\label{theo:cms}
Let $\cc{P} = \{\mathsf{P}_{\theta} : \theta \in \Theta\}$ be dominated by a $\sigma$--finite measure
$\nu$, with densities
\[
p_{\theta}(x) = \frac{d \mathsf{P}_{\theta}}{d \nu}(x).
\]
Define the support of the model by
\[
\spp(\cc{P}) = \{ x \in \cc{X} : \exists \theta \in \Theta \text{ such that } p_{\theta}(x) > 0 \}.
\]

Suppose that a statistic $T$ satisfies the following property: for all
$x, \tilde{x} \in \spp(\cc{P})$,
\begin{equation}
\label{eq:cms}
T(x) = T(\tilde{x})
\iff
\exists \, c(x,\tilde{x}) > 0 \text{ such that }
p_{\theta}(x) = p_{\theta}(\tilde{x}) \cdot c(x,\tilde{x})
\quad \forall \theta \in \Theta.
\end{equation}
Then $T$ is minimal sufficient for $\cc{P}$.
\end{theo}

\paragraph{Interpretation.}
The right--hand side of \eqref{eq:cms} states that the likelihood functions corresponding to
$x$ and $\tilde{x}$ are proportional as functions of $\theta$.
Thus, two data sets receive the same value of $T$ if and only if they induce the same likelihood
up to a multiplicative constant.
This criterion therefore identifies minimal sufficient statistics as those that classify
data sets exactly according to likelihood equivalence.

\paragraph{Comments}
\begin{itemize}
    \item Condition \eqref{eq:cms} provides a constructive way to identify minimal sufficient
    statistics by analyzing when likelihood functions are proportional.
    \item The criterion depends only on likelihoods and does not require computing conditional
    distributions.
    \item The original result is due to Lehmann and Scheff\'e (1950, Theorem~6.3).
\end{itemize}
\medskip
\begin{proof}[Discrete case]
We show that $T$ is sufficient and minimal sufficient.

\medskip
\noindent
\textbf{Step 1: $T$ is sufficient.}
For each value $t$ of $T$, choose a representative element
\[
x_t \in T^{-1}(\{t\}) = \{ x : T(x) = t \},
\]
with $x_t \in \spp(\cc{P})$ whenever
$T^{-1}(\{t\}) \cap \spp(\cc{P}) \neq \emptyset$.

For any $x \in \cc{X}$, we have $T(x) = T(x_{T(x)})$, so by assumption \eqref{eq:cms},
\[
p_{\theta}(x) = p_{\theta}(x_{T(x)}) \, c(x,x_{T(x)}),
\quad \theta \in \Theta,
\]
where we define $c(x,x_{T(x)}) = 0$ if $x \notin \spp(\cc{P})$.
This representation expresses $p_{\theta}(x)$ as a product of a term depending on $\theta$
only through $T(x)$ and a term independent of $\theta$.
Hence, by Neyman's factorization criterion, $T$ is sufficient.

\medskip
\noindent
\textbf{Step 2: $T$ is minimal sufficient.}
Let $T'$ be any other sufficient statistic and let
$x,\tilde{x} \in \spp(\cc{P})$ satisfy $T'(x) = T'(\tilde{x})$.
We must show that $T(x) = T(\tilde{x})$.

Since $T'$ is sufficient, the factorization criterion yields
\begin{align*}
p_{\theta}(x) &= g'_{\theta}(T'(x)) h'(x), \\
p_{\theta}(\tilde{x}) &= g'_{\theta}(T'(\tilde{x})) h'(\tilde{x}).
\end{align*}
Because $T'(x) = T'(\tilde{x})$, the $g'_{\theta}$ terms coincide, and thus
\begin{equation}
\label{eq:ch3_fin}
p_{\theta}(x) = p_{\theta}(\tilde{x}) \cdot \frac{h'(x)}{h'(\tilde{x})}
\quad \forall \theta \in \Theta.
\end{equation}
Since $\tilde{x} \in \spp(\cc{P})$, we have $h'(\tilde{x}) > 0$, so the ratio is well defined.
By condition \eqref{eq:cms}, this proportionality implies $T(x) = T(\tilde{x})$.

Therefore, $T$ is a function of every sufficient statistic $T'$, and hence $T$ is minimal
sufficient.
\end{proof}

\begin{ex}[Gaussian model]
Let $X_1,\dots,X_n$ be i.i.d.\ $\cc{N}(\mu,\sigma^2)$ with parameter
$\theta = (\mu,\sigma^2) \in \Theta = \R \times (0,\infty)$.
Recall that the maximum likelihood estimators are
\[
\hat{\mu} = \overline{X}_n,
\qquad
\hat{\sigma}^2 = \frac{1}{n}\sum_{i=1}^n (X_i - \overline{X}_n)^2 .
\]
\end{ex}

\begin{claim}
The statistic $(\hat{\mu},\hat{\sigma}^2)$ is minimal sufficient.
\end{claim}

\begin{proof}
The joint density of $\bs{X}=(X_1,\dots,X_n)$ is
\[
p_{\theta}(x)
=
\left( \frac{1}{\sqrt{2\pi\sigma^2}} \right)^n
\exp\!\left\{
-\frac{1}{2\sigma^2}\sum_{i=1}^n (x_i-\mu)^2
\right\}.
\]
Using the identity
\[
       \sum\limits_{i=1}^{n} (x_{i}-\mu)^2
       = \sum\limits_{i=1}^{n} (x_{i} - \bar{x}_{n}+\bar{x}_{n}-\mu)^2
       = \sum\limits_{i=1}^{n} (x_{i}-\bar{x}_{n})^2 + n(\bar{x}_{n}-\mu)^2,
\]
we obtain
\[
p_{\theta}(x)
\propto
\exp\!\left\{
-\frac{n}{2\sigma^2}
\bigl[(\bar{x}_n-\mu)^2 + \hat{\sigma}^2_x\bigr]
\right\}.
\]

Let $x,y \in \R^n$ be two data sets. The ratio of likelihoods is
\begin{align*}
\frac{p_{\theta}(x)}{p_{\theta}(y)}
&=
\exp\!\left\{
-\frac{n}{2\sigma^2}
\Bigl[
(\bar{x}_n-\mu)^2 - (\bar{y}_n-\mu)^2
+ \hat{\sigma}^2_x - \hat{\sigma}^2_y
\Bigr]
\right\} \\
&=
\exp\!\left\{
-\frac{n}{2\sigma^2}
\Bigl[
(\bar{x}_n^2-\bar{y}_n^2)
-2\mu(\bar{x}_n-\bar{y}_n)
+ \hat{\sigma}^2_x - \hat{\sigma}^2_y
\Bigr]
\right\}.
\end{align*}
This ratio is independent of $\theta=(\mu,\sigma^2)$ if and only if
\[
\bar{x}_n = \bar{y}_n
\quad\text{and}\quad
\hat{\sigma}^2_x = \hat{\sigma}^2_y .
\]
By the minimal sufficiency criterion of \autoref{sub:cms},
$(\hat{\mu},\hat{\sigma}^2)$ is minimal sufficient.
\end{proof}

\begin{ex}[Uniform location model]
Let $X_1,\dots,X_n$ be i.i.d.\ $\mathrm{Uniform}(\theta,\theta+1)$ for $\theta \in \R$.
Applying the minimal sufficiency criterion yields
\[
T(x) =
\Bigl(
\min\{x_1,\dots,x_n\},
\max\{x_1,\dots,x_n\}
\Bigr)
\]
as a minimal sufficient statistic; see Casella and Berger (2002,
Example~6.2.15) for details.

This example is noteworthy because the model is 1--dimensional,
yet the minimal sufficient statistic is 2--dimensional.
In fact, there are many low--dimensional models for which the order
statistics are minimal sufficient; see Lehmann and Casella (1998,
Example~1.6.15).
\end{ex}
